\section{Data}

%%% SUGGESTION: This section should probably become "Experimental Setup and Benchmark Families" or simply "Experiments" for conference readability.
%%% SUGGESTION: Right now the strongest missing piece is experimental discipline rather than additional ideas.
%%% SUGGESTION: Reviewers will mainly care about:
%%% SUGGESTION: - parameter-count parity
%%% SUGGESTION: - seed count
%%% SUGGESTION: - fair baselines
%%% SUGGESTION: - train/test methodology
%%% SUGGESTION: - whether transfer/interference metrics are rigorously defined

\section{Benchmark Families}

The full cartesian product of all datasets and all known activation patterns, loss functions, learning rate, optimizers now with the new rotation operator would be thorough, it's not quite feasible for the scope of a paper, but it is considered for future work and exhaustive Hexnet dynamic surveying.

Therefore we define a Benchmark Family is a subset of this cartesian product where hold certain hyperparameter constants and test individual themes with thess operators.
The theme of each family along with our findings are detailed below. The datasets are defined afterwards.

%%% SUGGESTION: Existing text: "all known activation patterns"
%%% SUGGESTION: Narrow this wording. Reviewers may interpret it literally.
%%% SUGGESTION: Existing text: "future work and exhaustive Hexnet dynamic surveying"
%%% SUGGESTION: Consider simpler language such as "future exhaustive hyperparameter sweeps".
%%% SUGGESTION: Define benchmark families formally.
%%% TODO: add a benchmark-family table summarizing:
%%% TODO: - family objective
%%% TODO: - datasets
%%% TODO: - fixed variables
%%% TODO: - swept variables
%%% TODO: - evaluation metrics

\subsection{Family A - Linear Structure Recovery}

%%% SUGGESTION: This is probably the correct first benchmark family.
%%% SUGGESTION: Linear and near-linear tasks are good because they isolate topology/scheduling behavior before nonlinear complexity dominates.

Hyper parameters that we will sweep:

\begin{table}
    \centering
    \begin{tabular}{c|c}
        Variable & Options \\
        \hline
        \hline
        Models & hex, mlp \\
        \hline
        Activations & linear, relu, leaky\_relu \\
        \hline
        Losses & mean\_squared\_error, huber \\
        \hline
        Learning Rates & constant, exponential\_decay \\
        \hline
        Dimension & Ladder up by twos, 2, 4, 6, 8 \\
    \end{tabular}
    \label{tab:Benchmark A Hyper parameters}
\end{table}

%%% TODO: add random seed count.
%%% TODO: add optimizer choices.
%%% TODO: add batch size.
%%% TODO: add train/validation/test split policy.
%%% TODO: add rotation schedule variables.
%%% TODO: explicitly state whether MLP baselines are parameter matched.
%%% SUGGESTION: Parameter-count parity may become the single most important reviewer concern.

\begin{table}
    \centering
    \begin{tabular}{c|c}
        Dataset & Notes \\
        \hline
        \hline
        identity & stuff \\
        \hline
        linear\_scale & stuff \\
        \hline
        diagonal\_scale & stuff \\
        \hline
        diagonal\_linear & stuff \\
        \hline
        full\_linear & stuff \\
        \hline
        low\_rank\_linear & stuff \\
        \hline
        affine & stuff \\
        \hline
        orthogonal\_rotation & stuff \\
        \hline
    \end{tabular}
\end{table}

%%% TODO: replace placeholder dataset notes.
%%% TODO: specify exact transformation equations.
%%% TODO: specify dimensionality.
%%% TODO: specify whether data are noisy or noiseless.
%%% TODO: specify sample counts.
%%% TODO: define orthogonal_rotation formally.

This helps us achieve parity with MLP with basic data sets and basic flow validation.
We see hex geometry perform just as well as equivalent MLP architecture, then it may be safe to assume isotropic maps are not trivially penalized.

We also see that on diagonal\_scale, diagonal\_linear, and low\_rank\_linear datasets, Hexnet performs better.

Hex Net loss and regression values fluctuate harshly after the rotation operator is applied, but this interference helps over all loss reach lower optimas on average.

When we specifically tested the orthogonal\_rotation dataset, we noticed that it differed in performance to the other affine transformations encoded in datasets.

%%% SUGGESTION: Move all empirical observations out of Data and into Results.
%%% SUGGESTION: The Data/Experimental Setup section should define methodology, not conclusions.
%%% SUGGESTION: Existing text: "Hexnet performs better"
%%% SUGGESTION: Add statistical framing later:
%%% SUGGESTION: - mean across seeds
%%% SUGGESTION: - variance
%%% SUGGESTION: - effect size
%%% SUGGESTION: - significance if appropriate
%%% SUGGESTION: Existing text: "fluctuate harshly"
%%% SUGGESTION: Replace with measurable language such as oscillatory convergence, transient loss spikes, or synchronization instability.

\subsection{Base Case}
Baselines for $H_0$ have been created where loss function, activation functions, and n are variable, and we can measure their performance with regression metrics such as scalar loss (according to the respective loss function), accuracy (root mean square deviation, as is standard for regression based data sets) and $R^2$ (the coefficient of determination).

%%% SUGGESTION: RMSE is not usually called "accuracy".
%%% SUGGESTION: Use separate regression metrics terminology.
%%% TODO: define all evaluation metrics formally.
%%% TODO: define transfer/interference metrics formally.
%%% TODO: define how cross-orientation evaluation is performed.

\subsection{Rotation Scheduling Families}

%%% TODO: add benchmark-family stubs for:
%%% TODO: - single-orientation training
%%% TODO: - sequential orientation rotation
%%% TODO: - train-one/test-all transfer
%%% TODO: - adjacent vs opposite orientation coupling
%%% TODO: - synchronization interval sweeps
%%% TODO: - EPR sweeps
%%% TODO: - rotation ordering sweeps

%%% SUGGESTION: This may become the most novel empirical contribution of the paper.
%%% SUGGESTION: The scheduling/interference behavior is likely more interesting than raw regression performance.

\subsection{Simulated Data Setup}

I did things here (identity dataset, linear dataset, chat gpt outputted)

%%% TODO: formally define synthetic data generation.
%%% TODO: specify random distributions.
%%% TODO: specify reproducibility seed policy.
%%% TODO: specify whether train/test distributions differ.
%%% SUGGESTION: Avoid mentioning "chat gpt outputted" in the final paper.

\subsection{Real Data Sets}

I did things with real data here

%%% TODO: define every real dataset.
%%% TODO: specify preprocessing.
%%% TODO: specify feature normalization.
%%% TODO: specify why equal-dimensional input/output assumptions remain valid.
%%% TODO: specify train/test split methodology.

\subsection{Baseline Definitions}

%%% TODO: add a baseline comparison table.
%%% TODO: include:
%%% TODO: - parameter count
%%% TODO: - topology
%%% TODO: - shared weights
%%% TODO: - number of subnetworks
%%% TODO: - schedule policy
%%% TODO: - inference procedure

%%% SUGGESTION: You probably do not need many baseline families for the first paper.
%%% SUGGESTION: Strong conference papers often win by having fewer but cleaner experiments.
%%% SUGGESTION: Recommended minimal baseline set:
%%% SUGGESTION: - parameter-matched MLP
%%% SUGGESTION: - HexNN single orientation
%%% SUGGESTION: - HexNN rotation-scheduled
%%% SUGGESTION: - six-independent-network ensemble baseline

\subsection{Experimental Controls}

%%% TODO: explicitly define:
%%% TODO: - parameter-count matching strategy
%%% TODO: - random seed count
%%% TODO: - stopping criteria
%%% TODO: - early stopping policy
%%% TODO: - hardware used
%%% TODO: - runtime budget
%%% TODO: - optimizer settings
%%% TODO: - gradient clipping policy if any
%%% TODO: - learning-rate scheduler policy

%%% SUGGESTION: These controls materially affect review outcomes.
%%% SUGGESTION: Reviewer confidence increases substantially when experiment controls are explicit.
